% Slide thuyết trình BTL số 2 – Social Golfer Problem bằng SAT
% Biên dịch bằng XeLaTeX (Overleaf: Compiler > XeLaTeX). Chỗ cần điền: \TODO{...}
\documentclass[aspectratio=169,11pt]{beamer}

\usepackage[vietnamese,english]{babel}
\usepackage{fontspec}
\setmainfont{texgyreheros}[
  Extension=.otf,
  UprightFont=*-regular,
  BoldFont=*-bold,
  ItalicFont=*-italic,
  BoldItalicFont=*-bolditalic]
\setmonofont{texgyrecursor}[
  Extension=.otf,
  UprightFont=*-regular,
  BoldFont=*-bold,
  ItalicFont=*-italic,
  BoldItalicFont=*-bolditalic]
\usepackage{amsmath,amssymb}
\usepackage{booktabs}
\usepackage{graphicx}
\usepackage{tikz}
\graphicspath{{figures/}}

% ---------- Giao diện ----------
\definecolor{navy}{HTML}{1E3A8A}
\definecolor{accent}{HTML}{2563EB}
\definecolor{ink}{HTML}{1E293B}
\usetheme{default}
\usecolortheme{default}
\setbeamercolor{structure}{fg=navy}
\setbeamercolor{normal text}{fg=ink}
\setbeamercolor{frametitle}{fg=white,bg=navy}
\setbeamercolor{title}{fg=navy}
\setbeamertemplate{navigation symbols}{}
\setbeamertemplate{itemize item}{\color{accent}$\blacktriangleright$}
\setbeamertemplate{footline}{%
  \hfill\insertframenumber/\inserttotalframenumber\hspace{1em}\vspace{0.6em}}
\setbeamerfont{frametitle}{series=\bfseries,size=\Large}
\addtobeamertemplate{frametitle}{}{\vspace{0.4em}}

\newcommand{\TODO}[1]{\textcolor{red}{[TODO: #1]}}
\newcommand{\RepositoryURL}{https://github.com/thnhmai06/SAT-SGP}

\title{Giải bài toán Social Golfer\\bằng SAT và so sánh với CP, ILP}
\subtitle{Bài tập lớn số 2 -- Các vấn đề hiện đại KHMT, HK I 2026--2027}
\author{Mai Hải Thành (24021627) \and Bùi Đình Cảnh (24021392)}
\institute{Trường Đại học Công nghệ, Đại học Quốc gia Hà Nội}
\date{\today}

\begin{document}

\begin{frame}[plain]
  \titlepage
\end{frame}

\begin{frame}{Nội dung}
  \tableofcontents
\end{frame}

\section{Bài toán}
\begin{frame}{Social Golfer Problem (SGP)}
  \begin{itemize}
    \item $n = g\cdot s$ người chơi, chia thành $g$ nhóm, mỗi nhóm $s$ người.
    \item Chơi $w$ tuần; mỗi tuần mỗi người thuộc đúng một nhóm.
    \item \textbf{Ràng buộc:} hai người bất kỳ cùng nhóm \emph{tối đa một lần}.
    \item Câu hỏi: có tồn tại lịch $w$ tuần hợp lệ cho $g$-$s$-$w$ không?
  \end{itemize}
  \vspace{0.5em}
  \TODO{Động cơ ứng dụng (lập lịch, thiết kế thí nghiệm) và độ khó.}
\end{frame}

\begin{frame}{Ví dụ: $g=3,\ s=3,\ w=4$ (9 người chơi)}
  \centering
  \begin{tabular}{cccc}
    \toprule
    Tuần 1 & Tuần 2 & Tuần 3 & Tuần 4 \\
    \midrule
    1 2 3 & 1 4 7 & 1 5 9 & 1 6 8 \\
    4 5 6 & 2 5 8 & 2 6 7 & 2 4 9 \\
    7 8 9 & 3 6 9 & 3 4 8 & 3 5 7 \\
    \bottomrule
  \end{tabular}

  \vspace{0.8em}
  Mỗi cặp người xuất hiện chung nhóm đúng một lần
  ($4\times 3\times 3 = 36 = \binom{9}{2}$ cặp).
\end{frame}

\section{Phương pháp}
\begin{frame}{Mã hóa SAT}
  \begin{itemize}
    \item Biến: $x_{p,r,k}$ -- người $p$ thuộc nhóm $k$ ở tuần $r$.
    \item Mỗi người đúng một nhóm mỗi tuần: ALO + AMO.
    \item Mỗi nhóm đúng $s$ người: ràng buộc cardinality.
    \item Mỗi cặp gặp nhau tối đa một lần: biến phụ ``cùng nhóm''.
    \item Phá đối xứng: \TODO{mô tả cách dùng}.
  \end{itemize}
\end{frame}

\begin{frame}{Các mã hóa AMO}
  \centering
  \begin{tabular}{lll}
    \toprule
    Mã hóa & Biến phụ & Mệnh đề \\
    \midrule
    Pairwise   & $0$ & $\binom{k}{2}$ \\
    Sequential & $\Theta(k)$ & $\Theta(k)$ \\
    Binary     & \TODO{} & \TODO{} \\
    Commander  & \TODO{} & \TODO{} \\
    Product    & \TODO{} & \TODO{} \\
    \bottomrule
  \end{tabular}
\end{frame}

\begin{frame}{Quy trình giải}
  \centering
  \begin{tikzpicture}[node distance=2.4cm,
      box/.style={draw=navy,rounded corners,thick,minimum width=2.3cm,minimum height=1cm,align=center,font=\small},
      arr/.style={->,thick,accent}]
    \node[box] (a) {Instance\\$(g,s,w)$};
    \node[box,right of=a,xshift=0.4cm] (b) {Encode\\CNF};
    \node[box,right of=b,xshift=0.4cm] (c) {SAT\\solver};
    \node[box,right of=c,xshift=0.4cm] (d) {Giải mã\\nghiệm};
    \node[box,right of=d,xshift=0.4cm] (e) {Checker\\độc lập};
    \draw[arr] (a)--(b); \draw[arr] (b)--(c); \draw[arr] (c)--(d); \draw[arr] (d)--(e);
  \end{tikzpicture}

  \vspace{1em}
  \TODO{SAT nhiều lần / Incremental SAT nếu có dùng.}
\end{frame}

\begin{frame}{Mô hình đối chứng}
  \begin{itemize}
    \item Cplex CP, Cplex MIP, Gurobi, CP-SAT (OR-Tools).
    \item Cùng tập lịch hợp lệ, cùng giới hạn thời gian.
    \item \TODO{Tham số solver, số luồng.}
  \end{itemize}
\end{frame}

\section{Thực nghiệm}
\begin{frame}{Thiết lập thực nghiệm}
  \begin{itemize}
    \item Instance: \TODO{CSPLib + instance lớn}.
    \item SAT solver: \TODO{CaDiCaL / Kissat / Glucose, qua PySAT hay trực tiếp}.
    \item Phần cứng: \TODO{CPU, RAM, OS}; giới hạn thời gian: \TODO{giây/ca}.
    \item Mọi nghiệm đều được kiểm tra bằng checker độc lập.
  \end{itemize}
\end{frame}

\begin{frame}{Kết quả: so sánh SAT encoding}
  % \includegraphics[width=0.85\linewidth]{encoding_time.pdf}  % vẽ bằng cnsplots
  \centering
  \TODO{Biểu đồ thời gian / số instance giải được theo từng encoding.}
\end{frame}

\begin{frame}{Kết quả: SAT so với CP và ILP}
  \centering
  \TODO{Bảng / biểu đồ SAT vs Cplex CP, Cplex MIP, Gurobi, CP-SAT.}
\end{frame}

\section{Kết luận}
\begin{frame}{Kết luận}
  \begin{itemize}
    \item \TODO{Kết quả chính (1--2 ý định lượng).}
    \item \TODO{Hạn chế.}
    \item \TODO{Hướng phát triển.}
  \end{itemize}
  \vspace{1em}
  Mã nguồn và kết quả: \url{\RepositoryURL}
\end{frame}

\begin{frame}[plain]
  \centering
  \vfill
  {\Huge\bfseries\color{navy} Cảm ơn thầy và các bạn đã lắng nghe!}\\[1em]
  Hỏi đáp
  \vfill
\end{frame}

\end{document}
